\begin{lstlisting}[language=c++,label={code:24_01},caption={一般線形モデルを実行するコード}]
dat <- read_csv("baseballDecade.csv") %>% filter(Year=="2020年度")
dataSet <- list(N = NROW(dat), Y = dat$height, X = dat$weight)
## cmdstanrで実行する場合
modelC <- cmdstan_model("LM.stan")
fit <- modelC$sample(
  data = dataSet,
  chains = 4,
  parallel_chains = 4,
  iter_warmup = 1000,
  iter_sampling = 2000
)
## 実行後のファイルをstanfit形式に置き換えておく
fit.stanfit <- fit$output_files() %>% rstan::read_stan_csv()
## 事後予測を取り出す
predY <- rstan::extract(fit.stanfit)$predY
# 事後予測分布の描画
bayesplot::ppc_dens_overlay(y = dataSet$Y, yrep = predY[1:10, ])
bayesplot::ppc_intervals(
    y = dataSet$Y,
    yrep = predY,
    x = dataSet$X,
    prob = 0.5,
    prob_outer = 0.95
)
\end{lstlisting}
